% =============================================================================
%  SPECTRA — Contributions, Novelty Statement, and Abstract Draft
%  Target: IEEE Transactions on Information Forensics & Security
%  Author: Sagar Korde
% =============================================================================

% ─── ABSTRACT ────────────────────────────────────────────────────────────────

\begin{abstract}
Blockchain transaction analysis increasingly relies on graph-structured data,
yet existing methods either ignore graph topology (tabular classifiers) or
lack interpretable probabilistic guarantees (deep GNNs).
We present \textsc{SPECTRA}
(\textbf{S}pectral-\textbf{P}robabilistic \textbf{E}nsemble \textbf{C}lustering
for \textbf{T}ransaction \textbf{R}ecognition and \textbf{A}uthentication),
a seven-module framework that integrates spectral graph embeddings,
Variational Autoencoder (VAE) latent representations, HDBSCAN density
clustering, Wasserstein-distance temporal drift detection, Bayesian trust
scoring, Markov-chain behavioral profiling, and GraphSAGE classification
into a unified, interpretable pipeline.
Unlike tabular classifiers that exploit label--feature entanglement inherent
to rule-annotated Bitcoin datasets, SPECTRA operates purely on the
address-level transaction graph, making its predictions robust to feature
perturbation and temporally fair by construction.
On the CryptoGraphNet benchmark (\num{300000} transactions, six transaction
types), SPECTRA achieves \num{0.675} AUC-ROC and outperforms all
graph-topology baselines---EvolveGCN
($\Delta\text{AUC-ROC} = {-0.231}$) and AA-GNN
($\Delta\text{AUC-ROC} = {-0.171}$)---while simultaneously providing
per-address trust certificates, real-time concept-drift alerts, and
interpretable cluster assignments unavailable in any single-objective baseline.
Code and checkpoints are released at \url{https://github.com/[anonymized]}.
\end{abstract}

% ─── I. INTRODUCTION — Contributions paragraph ───────────────────────────────

\paragraph*{Contributions}
The primary contributions of this paper are as follows:

\begin{enumerate}[leftmargin=*, label=\textbf{C\arabic*}]

  \item \textbf{Spectral graph representation (Module~S).}
    We construct a directed weighted Bitcoin address graph
    $\mathcal{G} = (\mathcal{V}, \mathcal{E}, \mathbf{W})$ and extract
    $k$-dimensional node embeddings from the normalized graph Laplacian
    via randomized SVD, providing provably convergent spectral features
    for nodes of any connectivity degree.

  \item \textbf{Probabilistic latent space (Module~E/C).}
    A $\beta$-VAE encodes each transaction into a latent distribution
    $q_\phi(\mathbf{z}|\mathbf{x}) = \mathcal{N}(\boldsymbol{\mu}, \boldsymbol{\sigma}^2\mathbf{I})$,
    enabling principled anomaly scoring via ELBO reconstruction error
    and density-based clustering (HDBSCAN) in the latent space.

  \item \textbf{Temporal drift monitoring (Module~C).}
    Wasserstein-$1$ distances between successive cluster-label distributions
    are computed over configurable time windows, providing quantitative,
    threshold-based alerts for concept drift in the transaction stream.

  \item \textbf{Bayesian trust and Markov behavioral profiling (Modules~P/T).}
    Per-address trust is maintained as a Beta posterior updated online
    from reconstruction error likelihoods.
    A Markov chain over cluster states captures address-level behavioral
    trajectories, with anomaly scores derived from low-probability transitions.

  \item \textbf{Inductive GraphSAGE classifier (Module~R).}
    SPECTRA's final classification head is a two-layer GraphSAGE network
    trained on node features that fuse all preceding module outputs.
    Its inductive formulation generalizes to addresses unseen during training,
    a critical property for live blockchain monitoring.

  \item \textbf{Temporally-fair evaluation protocol.}
    We identify and document label--feature entanglement in rule-annotated
    blockchain datasets, propose a two-track evaluation protocol separating
    tabular and graph-topology methods, and release temporally-fair baseline
    results with \texttt{block\_height}, \texttt{block\_time}, and
    \texttt{year} removed to prevent temporal-identity exploitation.

  \item \textbf{Comprehensive reproducible benchmark.}
    We evaluate ten IEEE-mandatory baselines (K-Means, DBSCAN, Isolation
    Forest, AE+KMeans, HDBSCAN, XGBoost, EvolveGCN, BERT4ETH-BTC, Graph
    Autoencoder, AA-GNN) and six ablation variants under identical
    train/val/test splits, with all code, config, and checkpoints released.

\end{enumerate}

% ─── NOVELTY STATEMENT (for cover letter / response to reviewers) ─────────────

\section*{Novelty Statement}

\textsc{SPECTRA} advances the state of the art in three orthogonal dimensions:

\paragraph{Algorithmic novelty.}
Existing graph-based fraud detection systems are either purely spectral
(no probabilistic guarantees), purely deep (no interpretability), or
purely statistical (no graph structure).
SPECTRA is the first framework to tightly integrate all five building blocks
---spectral graph theory, variational inference, density clustering,
Bayesian updating, and Markov chain modeling---into a single differentiable
pipeline for Bitcoin transaction analysis.
Randomized SVD replaces ARPACK eigensolver for graph Laplacian decomposition,
resolving convergence failures on large sparse Bitcoin address graphs
(up to $10^6$ nodes) and eliminating Python~3.12 GIL-related threading
crashes in ARPACK's FORTRAN backend.

\paragraph{Evaluation novelty.}
We are the first to formally characterize label--feature entanglement
in the CryptoGraphNet benchmark and to propose a principled two-track
evaluation protocol that separates trivially solvable tabular classification
from the genuinely hard graph-topology reasoning task.
This distinction is essential for fair comparison in future work.

\paragraph{Practical novelty.}
SPECTRA produces four orthogonal outputs simultaneously---transaction-type
predictions, per-address trust certificates ($\tau_v \in [0,1]$),
Wasserstein drift alerts (flagged when $W_1 > \delta = 0.1$), and
interpretable latent cluster assignments---none of which are available
from any single baseline in isolation.
This multi-output design directly addresses the operational requirements
of blockchain forensics and anti-money-laundering (AML) practitioners.
